\section{Conclusion}\label{sec:conclusion}

A quantum slice helps a developer only if it never hides code that can change the qubits the developer asks about.
This paper defined that requirement as preservation of the qubits' joint state for every input together with
completeness for semantic dependence, and showed that one rule meets it, the write-all rule, under which a gate writes
every qubit it touches. Its slice is the backward light cone extended along classical bits, it needs only classical
data, classical control and quantum data dependence, and it is computed in one pass. A pre-registered study measured
what every other rule loses. Reading the control of a gate as unchanged gives wrong slices that often reproduce the
program's outcomes, entanglement edges add only code that cannot matter, and pruning trades dependences for size and
holds for one observable or one input. The QDG slice comes within a few points of the smallest slice that keeps every
dependence and costs milliseconds.

Two questions remain. The first is whether these slices shorten the manual reading that developers report does not
scale, which needs a study with developers. The second is how to compute them on the source developers write, statically and at the
scale of real programs, including loops that run on the device.
